%% kapa: exam=1 round=27 year=2016 period=1 form=A subject=economics subject_ko=경제학원론 seq=8 number=48
%% source: Q-Net 감정평가사 기출문제 (cst00309) artlSeq=5169402 file=extracted/1008/5169402/제27회 감정평가사 1차 문제/제27회 감정평가사 1교시문제(A형).pdf
%% conversion: kordoc 4.24.51 -> tools/gen1.py (auto); review=proofread-sonnet; flags=none
\begin{question}{48}
甲기업의 단기 총비용함수가 $C = 25 + 5Q$일 때, 甲기업의 단기 비용에 관한 설명으로 옳은 것은? (단, $Q$는 양(+)의 생산량이다.)
\begin{choices}
  \item 모든 생산량 수준에서 평균 가변비용과 한계비용은 같다.
  \item 모든 생산량 수준에서 평균 고정비용은 일정하다.
  \item 생산량이 증가함에 따라 한계비용은 증가한다.
  \item 평균비용 곡선은 U자 형태이다.
  \item 생산량이 일정 수준 이상에서 한계비용이 평균비용을 초과한다.
\end{choices}
\end{question}
\answer{1}
